\documentclass[10pt]{article}
\usepackage[T1]{fontenc}
\usepackage[utf8]{inputenc}
\usepackage[french]{babel}
\usepackage{amsmath,amssymb,mathrsfs,mathpazo}
\usepackage{geometry}
\geometry{a4paper,margin=2cm}
\usepackage{graphicx,array,tabularx,tabularray,colortbl,multirow,multicol,booktabs,enumitem,pifont,ulem,fancybox,diagbox,cancel}
\usepackage{pstricks,pstricks-add,pst-plot,pst-tree,pst-node,pst-text,pst-func,pst-math,pst-bspline,pst-3dplot,pst-eucl,pst-solides3d}
\let\pscalculate\undefined
\usepackage{pst-bezier}
\usepackage{tikz,pgfplots,tkz-tab}
\usetikzlibrary{arrows,arrows.meta,patterns,calc,positioning,trees,decorations.markings,shapes.geometric,shapes.misc,shapes.arrows,intersections,decorations.pathmorphing,decorations.pathreplacing,backgrounds,shadows,babel}
\pgfplotsset{compat=1.18}
\usepgfplotslibrary{fillbetween}
\usepackage[np]{numprint}
\usepackage{listings,makecell,pas-tableur,twoopt}
\let\barre\undefined
\usepackage[pstricks]{tablvar}
\definecolor{darkgreen}{RGB}{0,100,0}
\providecommand{\red}{\color{red}}
\providecommand{\blue}{\color{blue}}
\providecommand{\green}{\color{green}}
\providecommand{\black}{\color{black}}
\providecommand{\white}{\color{white}}
\providecommand{\gray}{\color{gray}}
\providecommand{\cyan}{\color{cyan}}
\providecommand{\magenta}{\color{magenta}}
\providecommand{\yellow}{\color{yellow}}
\providecommand{\violet}{\color{violet}}
\providecommand{\orange}{\color{orange}}
\providecommand{\euro}{\texteuro}
\providecommand{\EUR}{\texteuro}
\providecommand{\ieme}{\textsuperscript{e}}
\providecommand{\degres}{\ensuremath{^\circ}}
\providecommand{\degree}{\ensuremath{^\circ}}
\providecommand{\vv}[1]{\overrightarrow{#1}}
\providecommand{\Coord}[1]{#1}
\providecommand{\decosix}{}
\providecommand{\starredbullet}{\ensuremath{\bullet}}
\pagestyle{empty}
\setlength{\parindent}{0pt}
\usepackage{adjustbox}\DeclareUnicodeCharacter{2236}{:}\DeclareUnicodeCharacter{2190}{\ensuremath{\leftarrow}}
\providecommand{\R}{\mathbb{R}}
\providecommand{\N}{\mathbb{N}}
\providecommand{\Z}{\mathbb{Z}}
\providecommand{\Q}{\mathbb{Q}}
\providecommand{\D}{\mathbb{D}}
\providecommand{\C}{\mathbb{C}}
\providecommand{\e}{{\mathrm{e}}}
\providecommand{\vect}[1]{\overrightarrow{#1}}
\providecommand{\vectt}[1]{\overrightarrow{\mathrm{#1}}}
\providecommand{\vvt}[1]{\overrightarrow{\mathrm{#1}}}
\providecommand{\np}[2][]{{#2~#1}}
\providecommand{\vv}[1]{{\overrightarrow{#1}}}
\renewcommand{\d}{{\mathrm{d}}}
\providecommand{\degres}{{^\circ}}
\providecommand{\degre}{{^\circ}}
\providecommand{\blue}{\color{blue}}
\providecommand{\red}{\color{red}}
\providecommand{\enskip}{\quad}
\providecommand{\llbracket}{\lbrack\!\lbrack}
\providecommand{\rrbracket}{\rbrack\!\rbrack}
\providecommand{\euro}{{\text{€}}}
\providecommand{\no}{{\text{n°}}}
\providecommand{\interval}[2]{{[#1;#2]}}
\providecommand{\intervFO}[2]{{[#1;#2[}}
\providecommand{\dotfill}{\ldots}

\providecommand{\ds}{\displaystyle}
\providecommand{\textding}[1]{#1}
\providecommand{\Oij}{\(\mathrm{O};\vec{\imath},\vec{\jmath}\)}
\providecommand{\Oijk}{\(\mathrm{O};\vec{\imath},\vec{\jmath},\vec{k}\)}
\providecommand{\Ouv}{\(\mathrm{O};\vec{u},\vec{v}\)}
\providecommand{\intervFF}[2]{\ensuremath{[#1;#2]}}
\renewcommand{\euro}{\ifmmode\text{\texteuro}\else\texteuro\fi}
\DeclareUnicodeCharacter{00B0}{\ifmmode^{\circ}\else\textdegree\fi}
\DeclareUnicodeCharacter{20AC}{\ifmmode\text{\texteuro}\else\texteuro\fi}
% Énoncés : Éducation nationale. Transcription et corrigés : APMEP, Denis Vergès et contributeurs. Découpage et classement : CoopMaths, Sébastien Lozano, Éric Elter et contributeurs. Adaptation de la présentation : Maths Explorer.
% CC BY-SA : https://coopmaths.fr/www/about/
% https://coopmaths.fr/alea/static/e3c/2020/tex/e3c_2020_00_sujet55_4.tex
\begin{document}
\noindent\textit{APMEP / CoopMaths. CC BY-SA. Présentation adaptée par Maths Explorer.}\par\medskip
\medskip

La bibliothèque municipale étant devenue trop petite, une commune a décidé d'ouvrir une médiathèque qui pourra contenir \np{100000} ouvrages au total. Pour l'ouverture prévue le 1\up{er} janvier 2020, la médiathèque dispose du stock de \np{35000} ouvrages de l'ancienne bibliothèque, augmenté de \np{7000} ouvrages supplémentaires neufs offerts par la commune.

\medskip

\textbf{Partie A}

\medskip

Chaque année, le bibliothécaire est chargée de supprimer 5\,\% des ouvrages, trop vieux ou abîmés, et d'acheter \np{6 000} ouvrages neufs.

On appelle $u_n$ le nombre, en milliers, d'ouvrages disponibles le 1\up{er} janvier de l'année (2020+$n$).

On donne $u_0 = 42$.

\medskip
 
\begin{enumerate}
\item  Justifier que, pour tout entier naturel n, on a $u_{n+1} = u_n\times  0,95+6$.
\item On propose ci-dessous un programme en langage Python :
\begin{lstlisting}[style=python]
def suite(n) :
u=42
for i in range(n) :
	u=0.95*u+6
return u
\end{lstlisting}

Expliquer ce que permet de déterminer ce programme.
\end{enumerate}

\medskip

\textbf{Partie B}

\medskip

La commune doit finalement revoir ses dépenses à la baisse, elle ne pourra financer que \np{4000} nouveaux ouvrages par an au lieu des \np{6000} prévus.

On appelle $v_n$ le nombre, en milliers, d'ouvrages disponibles le 1\up{er} janvier de l'année (2020+$n$).

\medskip

\begin{enumerate}
\item  On admet que $v_{n+1} = 0,95\times v_n + 4$ pour tout entier naturel $n\geqslant 0$ avec $v_0 = 42$.

On considère la suite $(w_n)$ définie, pour tout entier naturel $n$, par $w_n= v_n-80$.
	\begin{enumerate}
		\item Montrer que $\left(w_n\right)$ est une suite géométrique de raison $q = 0,95$ et préciser son premier terme $w_0$.
		\item En déduire l'expression de $w_n$ puis de $ v_n$ en fonction de $n$.

	\end{enumerate}
\item On donne ci-dessous un programme en langage Python.

\begin{center}
\begin{lstlisting}[style=python]
def objet(A) :
v=42
n=0
while v<A :
	v=0.95*v+4
	n=n+1
return n 
\end{lstlisting}
\end{center}

L'appel à la fonction objet(70) renvoie 27. Interpréter ce résultat dans le contexte de l'exercice.
\end{enumerate}

\end{document}
